El Quetzal-2 es un nanosatélite CubeSat 2U desarrollado como proyecto académico por estudiantes y \textit{staff} de la Universidad del Valle de Guatemala el cual pone a prueba una \textit{On-Board Computer} (OBC) diseñada localmente que ejecuta una inteligencia artificial capaz de identificar nubes. Este nanosatélite cuenta con una computadora a bordo compuesta de por dos unidades de procesamiento, una OBC principal basada en hardware comercial de GomSpace y otra OBC secundaria desarrollada localmente utilizando Portentas H7 de Arduino.

Las OBC son sistemas vitales dentro de una misión espacial puesto que gestionan todos los subsistemas del satélite así como la comunicación entre estos. Este proyecto se centra en las bases de esta OBC local a nivel de \textit{software}, sentando las bases para los distintos modos de operación que necesita la misión, la toma de fotografías y la interoperabilidad entre las dos unidades de procesamiento que componen la OBC completa.

El impacto de esta OBC desarrollada localmente representa un gran punto de inicio para próximas misiones espaciales, bajando aún más la barrera económica de entrada para desarrollar misiones espaciales con respecto a soluciones \textit{of-the-shelf}, puesto que demuestra que no es necesario el uso de \textit{hardware} altamente costoso con muchas más protecciones para llevar a cabo una misión exitosa.
